\newpage 
\appendix
\section*{Appendix}



\section{Statistical Methodology}\label{apx:stats_results}

In Table \ref{tab:all_results}, we report the mean and an error bar for each reported method on the three benchmarks. To obtain the error bar we used a simple Wald 95\% confidence interval for the proportion. The Wald confidence interval assumes a normal approximation for the sample mean which is only valid for larger sample sizes and is only accurate for proportions not near 0 or 1. Our application meets these criteria: the smallest evaluated test set consists of 181 samples, and all reported results hover around 30\% with the exception of GPT-4. We also computed confidence intervals using the Wilson interval and found similar results (though Wilson intervals are not symmetric). For simplicity, we report only Wald intervals here. 

We also report the results of a statistical test comparing \team (GPT-4o, o1) to each reported method. We used the z-test to compare the accuracy of \team to each baseline in Tabe \ref{tab:all_results}. The z-test for proportions is the only feasible test in our setting because we only have the mean accuracy and not results on each example (the task-level test set results are hidden by the benchamrks). Therefore, we cannot apply McNemar's Test or a pairwise t-test. The limitation of the z-test is that it ignores pairing of the data and is generally conservative. We hope future benchmark leaderboards can release confidence intervals in addition to the reported mean.

\section{Capability to Category Mapping}\label{apx:capability_mapping}

Figure \ref{fig:ablations} (right) shows task performance results broken down by capabilities required. These capabilities are based on the capabilities human annotators reported as needed to solve tasks in the GAIA benchmark dataset \cite{mialon2023gaiabenchmarkgeneralai}. We organized and re-coded these annotations into the following categories to better reflect the roles of \team's agents:

\begin{itemize}
    \item \textbf{Web browsing:} capabilities related to searching and browsing the web. Examples: \texttt{web browser}, \texttt{search engine}, \texttt{maps}, \texttt{access to internet archives}
    \item \textbf{Coding:} capabilities related to coding and execution. Examples: \texttt{Coding}, \texttt{Python}, \texttt{calculator}, \texttt{audio/video processing}, \texttt{text processing}, \texttt{natural language processing}
    \item \textbf{File handling:} capabilities related to handling diverse file types. Examples: \texttt{Pdf viewer}, \texttt{Word}, \texttt{Excel}, \texttt{Powerpoint file access}, \texttt{CSV file access}, \texttt{XML file access}
    \item \textbf{No tools:} capabilities that can be performed inherently by agents, without tool-use, using multi-modal models. Examples: \texttt{Image recognition}, \texttt{OCR}, \texttt{Computer vision}, \texttt{color recognition}, \texttt{extracting text from images}
\end{itemize}

\newpage
\section{Error Analysis Code Book}
\label{sec:appendix:error}

In this section, we provide the final codes from the automated analysis of \team's behavior in the validation logs across all benchmarks. The codes are sorted by how often they appeared in the samples they were taken from. For each code, we include a definition and summaries of examples from the logs that were assigned that code.

{\small
\begin{longtable}{p{0.2\textwidth} p{0.3\textwidth} p{0.37\textwidth}}
\hline
\textbf{Name} & \textbf{Definition} & \textbf{Examples} \\
\hline
persistent-inefficient-actions &
Agents engaged in unproductive patterns without adapting despite facing failures. Ineffective strategies persisted, leading to delays and insufficient task outcomes. &
\begin{itemize}
  \item Agents clicked the same webpage sections multiple times, achieving no advancement in information retrieval.
  \item Agents unnecessarily engaged in general web searches instead of focusing on specified database tools.
  \item WebSurfer continued failing searches with no query modification, ignoring unsuccessful outcomes.
  \item The orchestrator commanded agents to use a flawed path repeatedly, indicating ongoing cycles without modification.
  \item Unaltered processes accessed incorrect data sets, consuming resources without advancing task goals.
\end{itemize}
\\ \hline

insufficient-verification-steps &
Tasks were marked complete without thorough data validation, leading to unreliable outcomes. Essential checks were skipped, resulting in erroneous assumptions about data integrity. &
\begin{itemize}
  \item Final outputs were accepted without validating data correctness, leading to potential errors.
  \item An orchestrator concluded a task although necessary criteria were unverified, risking incomplete achievement.
  \item A dataset's verification steps were skipped, causing unaddressed errors in downstream analysis.
  \item Document scans lacked quality checks before storing, leading to unreliable information in reports.
  \item Data interpretation lacked cross-verification assurances, uncovering gaps in computation assurances.
\end{itemize}
\\ \hline

underutilized-resource-options &
Agents consistently did not utilize available data, tools, or resources effectively. This resulted in inefficient task execution and repeated manual actions, even when automation was an option. &
\begin{itemize}
  \item Agents failed to integrate accessible descriptions, opting for redundant manual inputs despite metadata availability.
  \item Manual downloads persisted when FileSurfer was available for rapid document retrieval, unnecessarily extending the task.
  \item Available APIs were bypassed in favor of manual approaches, extending the task duration unnecessarily.
  \item Advanced search functions were overlooked, leading to reliance on broad manual explorations, slowing down the process.
  \item Complex data visualization was attempted with basic charting tools instead of comprehensive graphing tools.
\end{itemize}
\\ \hline

inefficient-navigation-attempts &
Errors occurred because of incorrect or inefficient navigation, leading to missed targets or prolonged task completion. Agents misinterpreted interface layouts and took inefficient paths. &
\begin{itemize}
  \item The agent mistakenly cycled through multiple tabs to find the 'Settings' page, resulting in delayed task progress.
  \item Incorrectly clicked navigation bars led to a user failing to access the proper configuration settings.
  \item Confusion over the UI design led an agent back to the main menu instead of the subsection needed for task completion.
  \item Cycling through history tabs resulted in missed current transaction logs.
  \item The agent persistently accessed the wrong page links, causing delays in retrieving important data.
\end{itemize}
\\ \hline

ineffective-team-communication &
Agents did not communicate information or direction effectively, causing task overlap and confusion. Miscommunications led to redundant work and uncoordinated actions. &
\begin{itemize}
  \item Two agents simultaneously accessed order histories due to unclear task division, resulting in duplicated efforts.
  \item Agents struggled to understand task requirements due to incomplete directive descriptions.
  \item The WebSurfer and FileSurfer failed to communicate their findings, leaving the Orchestrator to make decisions based on incomplete information.
  \item Repeated task restarts occurred because instructions were too vague, leading to incorrect interpretations.
  \item Agents observed unclear role demarcations that caused them to erroneously overwrite each other's progress.
\end{itemize}
\\ \hline

neglected-error-notifications &
Agents ignored known errors or warnings, allowing issues to persist and recur. This resulted in repeated inefficiencies and bottlenecks in task execution. &
\begin{itemize}
  \item Repeated 'ValueError' messages were ignored, allowing the underlying issue to continue unaddressed.
  \item Agents accessed 'hotel booking' pages multiple times despite 'No results found' errors repeatedly visible.
  \item System errors were logged without corrective attempts, leading to prolonged delays in task completion.
  \item Unresolved input errors led to repeated submissions of the same query form without variations.
  \item Despite validation warnings, agents continued with deprecated approaches, leading to unresolved problems.
\end{itemize}
\\ \hline

flawed-technical-implementations &
Tasks faced issues due to incorrect application of technical logic or processes. Misapplied techniques led to errors and inefficiencies. &
\begin{itemize}
  \item Agents encountered syntax errors in scripts due to incorrect indentation, halting task progress.
  \item Repeated runtime errors occurred as agents submitted misformatted queries without cross-verifying syntax.
  \item In accessing web modules, agents misapplied parameter mappings, causing non-responsive functions.
  \item Misinterpretation of calculation steps led to varying billing inconsistencies within an ecommerce task.
  \item Agents faced challenges in aligning sales data, resulting in flawed revenue projections.
\end{itemize}
\\ \hline

access-and-security-barriers &
Tasks were hindered due to security or access restrictions, affecting data integrity. Agents struggled with authentication, allowing unauthorized access or limiting task completion. &
\begin{itemize}
  \item Visible password fields were used, risking data exposure to unauthorized parties.
  \item Repeated attempts to submit forms were met with unresolved errors blocking access to data queries.
  \item Token visibility in the logs suggested possible unauthorized data access due to weak credential management.
  \item Insufficient password policies allowed repeated login attempts without user notifications, risking security.
  \item Repeated security filter alerts appeared as agents attempted to access restricted areas.
\end{itemize}
\\ \hline

delayed-feedback-updates &
There was a delay in communicating task progress or results, causing confusion and hindering coordination. Timeliness in updates was lacking, leaving tasks in ambiguity. &
\begin{itemize}
  \item After WebSurfer confirmed a task, the orchestrator delayed updating the user, leading to confusion.
  \item Notification of successful strategy analysis reached the user well after execution, causing temporary uncertainty.
  \item Despite quick report generation, notifications lagged behind, leading users to question task progression.
  \item Task completions were assumed due to lagged updates in user notifications post-webpage completions.
  \item Feedback on document completions took significantly longer to communicate to concerned tasks than average.
\end{itemize}
\\ \hline

successful-task-completions &
Tasks were completed smoothly with no errors, meeting all objectives efficiently. Agents achieved success through coordinated efforts, correct usage of resources, and thorough verification. &
\begin{itemize}
  \item During the data migration task, all user details were accurately uploaded and verified without issues.
  \item Correct transaction processing led to flawless financial reconciliation, with all figures matching expectations.
  \item Backend updates improved processing speeds, which performance tests later verified as positive and expected.
  \item Users experienced smooth account creations, following flawless processes without discrepancies or errors.
  \item Revisions in manuscript drafts met all editorial instructions perfectly, with edits applied and aligned as needed.
\end{itemize}
\\ \hline

\end{longtable}

}
